\documentclass[10pt]{article}
\usepackage[T1]{fontenc}
\usepackage[utf8]{inputenc}
\usepackage[french]{babel}
\usepackage{amsmath,amssymb,mathrsfs,mathpazo}
\usepackage{geometry}
\geometry{a4paper,margin=2cm}
\usepackage{graphicx,array,tabularx,tabularray,colortbl,multirow,multicol,booktabs,enumitem,pifont,ulem,fancybox,diagbox,cancel}
\usepackage{pstricks,pstricks-add,pst-plot,pst-tree,pst-node,pst-text,pst-func,pst-math,pst-bspline,pst-3dplot,pst-eucl,pst-solides3d}
\let\pscalculate\undefined
\usepackage{pst-bezier}
\usepackage{tikz,pgfplots,tkz-tab}
\usetikzlibrary{arrows,arrows.meta,patterns,calc,positioning,trees,decorations.markings,shapes.geometric,shapes.misc,shapes.arrows,intersections,decorations.pathmorphing,decorations.pathreplacing,backgrounds,shadows,babel}
\pgfplotsset{compat=1.18}
\usepgfplotslibrary{fillbetween}
\usepackage[np]{numprint}
\usepackage{listings,makecell,pas-tableur,twoopt}
\let\barre\undefined
\usepackage[pstricks]{tablvar}
\definecolor{darkgreen}{RGB}{0,100,0}
\providecommand{\red}{\color{red}}
\providecommand{\blue}{\color{blue}}
\providecommand{\green}{\color{green}}
\providecommand{\black}{\color{black}}
\providecommand{\white}{\color{white}}
\providecommand{\gray}{\color{gray}}
\providecommand{\cyan}{\color{cyan}}
\providecommand{\magenta}{\color{magenta}}
\providecommand{\yellow}{\color{yellow}}
\providecommand{\violet}{\color{violet}}
\providecommand{\orange}{\color{orange}}
\providecommand{\euro}{\texteuro}
\providecommand{\EUR}{\texteuro}
\providecommand{\ieme}{\textsuperscript{e}}
\providecommand{\degres}{\ensuremath{^\circ}}
\providecommand{\degree}{\ensuremath{^\circ}}
\providecommand{\vv}[1]{\overrightarrow{#1}}
\providecommand{\Coord}[1]{#1}
\providecommand{\decosix}{}
\providecommand{\starredbullet}{\ensuremath{\bullet}}
\pagestyle{empty}
\setlength{\parindent}{0pt}
\usepackage{adjustbox}\DeclareUnicodeCharacter{2236}{:}\DeclareUnicodeCharacter{2190}{\ensuremath{\leftarrow}}
\providecommand{\R}{\mathbb{R}}
\providecommand{\N}{\mathbb{N}}
\providecommand{\Z}{\mathbb{Z}}
\providecommand{\Q}{\mathbb{Q}}
\providecommand{\D}{\mathbb{D}}
\providecommand{\C}{\mathbb{C}}
\providecommand{\e}{{\mathrm{e}}}
\providecommand{\vect}[1]{\overrightarrow{#1}}
\providecommand{\vectt}[1]{\overrightarrow{\mathrm{#1}}}
\providecommand{\vvt}[1]{\overrightarrow{\mathrm{#1}}}
\providecommand{\np}[2][]{{#2~#1}}
\providecommand{\vv}[1]{{\overrightarrow{#1}}}
\renewcommand{\d}{{\mathrm{d}}}
\providecommand{\degres}{{^\circ}}
\providecommand{\degre}{{^\circ}}
\providecommand{\blue}{\color{blue}}
\providecommand{\red}{\color{red}}
\providecommand{\enskip}{\quad}
\providecommand{\llbracket}{\lbrack\!\lbrack}
\providecommand{\rrbracket}{\rbrack\!\rbrack}
\providecommand{\euro}{{\text{€}}}
\providecommand{\no}{{\text{n°}}}
\providecommand{\interval}[2]{{[#1;#2]}}
\providecommand{\intervFO}[2]{{[#1;#2[}}
\providecommand{\dotfill}{\ldots}

\providecommand{\ds}{\displaystyle}
\providecommand{\textding}[1]{#1}
\providecommand{\Oij}{\(\mathrm{O};\vec{\imath},\vec{\jmath}\)}
\providecommand{\Oijk}{\(\mathrm{O};\vec{\imath},\vec{\jmath},\vec{k}\)}
\providecommand{\Ouv}{\(\mathrm{O};\vec{u},\vec{v}\)}
\providecommand{\intervFF}[2]{\ensuremath{[#1;#2]}}
\renewcommand{\euro}{\ifmmode\text{\texteuro}\else\texteuro\fi}
\DeclareUnicodeCharacter{00B0}{\ifmmode^{\circ}\else\textdegree\fi}
\DeclareUnicodeCharacter{20AC}{\ifmmode\text{\texteuro}\else\texteuro\fi}
% Énoncés : Éducation nationale. Transcription et corrigés : APMEP, Denis Vergès et contributeurs. Découpage et classement : CoopMaths, Sébastien Lozano, Éric Elter et contributeurs. Adaptation de la présentation : Maths Explorer.
% CC BY-SA : https://coopmaths.fr/www/about/
% https://coopmaths.fr/alea/static/bac/2024/tex/bac_2024_05_sujet1_ameriquenord_2_cor.tex
\begin{document}
\noindent\textit{APMEP / CoopMaths. CC BY-SA. Présentation adaptée par Maths Explorer.}\par\medskip
\medskip















\begin{enumerate}
\item On considère les points A(1~;~0~;~3) et B(4~;~1~;~0).

Une représentation paramétrique de la droite (AB) est:























On a $\vect{\text{AB}}\begin{pmatrix}3\\1\\-3\end{pmatrix}$.
$M(x~;~y~;~z) \in (\text{AB}) \iff \vect{\text{A}M} = t\vect{\text{AB}}, \: t \in \R \iff $

$\left\{\begin{array}{l c l}x - 1&=&3t\\y - 0&=&1t\\z - 3&=&- 3t, \end{array}\right. \: t \in \R\iff \left\{\begin{array}{l c l}
x &=& 1 + 3t\\
y &=&\phantom{1 + 3}t\\
z&=&3 - 3t
\end{array}\right.\: t \in \R$ Réponse \textbf{c.~}

On considère la droite $(d)$ de représentation paramétrique 

$\left\{\begin{array}{l c l}
x&=&3 + 4t\\
y&=&\phantom{3 + }6t\\
z&=&4 - 2 t
\end{array}\right.$ avec $t \in \R$

\item 





Réponse \textbf{d}.
\item 













$\bullet~~$Les deux droites ont pour vecteurs directeurs respectifs : $\begin{pmatrix}4\\6\\-2\end{pmatrix}$ et $\begin{pmatrix}3\\-2\\1\end{pmatrix}$ : ces vecteurs ne sont pas colinéaires, les droites ne sont pas parallèles.

$\bullet~~$Les deux droites sont sécantes s'il existe deux réels $t$ et $k$ tels que :

$\left\{\begin{array}{l c l}
3 + 4t&=&- 2 + 3k\\
6t&=&- 1 - 2k\\
4 - 2t&=&1 + k
\end{array}\right. \iff \left\{\begin{array}{l c l}
4t&=&3k - 5\\6t&=&- 2k - 1\\
- 2t&=&k - 3\end{array}\right.$

Par différence de la ligne 1 et de la ligne 2, on obtient$ -2t = 5k - 4 = k - 3$ (ligne 3) soit $4k = 1 \iff k = \dfrac14$, puis comme $- 2t = k - 3 = \dfrac14 - 3 = - \dfrac{11}{4}$, d'où $t = \dfrac{11}{8}$, mais la première ligne donne $4t = 3k - 5 = \dfrac34 - 5 = - \dfrac{17}{4}$, d'où $t = \dfrac{17}{8}$ : ceci n'est pas possible : les deux droites ne sont pas sécantes.

$\bullet~~$Les deux droites ne sont pas confondues : il reste la réponse \textbf{b.}

\item  







Le vecteur $\vect{n}$ vecteur directeur de $(d)$ est un vecteur normal au plan $(P)$ . On a donc 

$M(x~;~y~;~z) \in (P) \iff \vect{\text{I}M} \cdot \vect{n} = 0 \iff 4(x - 2) + 6(y - 1) - 2(z - 0) = 0 \iff 4x - 8 + 6y - 6 - 2z = 0 \iff 4x + 6y -2z - 14 = 0 \iff 2x + 3y - z - 7 = 0$ : réponse \textbf{a.}
\end{enumerate}

\bigskip

\end{document}
